\textbf{1. Kiến thức cần nhớ}
\begin{itemize}
\item \textbf{Định nghĩa:} Cho mặt phẳng $(\alpha)$ và đường thẳng $\Delta$ cắt $(\alpha)$. Với mỗi điểm $M$ trong không gian, gọi $M'$ là giao điểm của $(\alpha)$ với đường thẳng đi qua $M$ và song song (hoặc trùng) với $\Delta$. Phép đặt tương ứng mỗi điểm $M$ với điểm $M'$ gọi là \textit{phép chiếu song song} lên mặt phẳng $(\alpha)$ theo phương $\Delta$. Trong đó: $(\alpha)$ là mặt phẳng chiếu, $\Delta$ là phương chiếu, $M'$ là hình chiếu song song của $M$.
\item \textbf{Tính chất} (với các đường thẳng không song song hoặc trùng với phương chiếu $\Delta$):
  \begin{itemize}
  \item Biến ba điểm thẳng hàng thành ba điểm thẳng hàng và giữ nguyên thứ tự của chúng.
  \item Biến đường thẳng thành đường thẳng, tia thành tia, đoạn thẳng thành đoạn thẳng.
  \item Biến hai đường thẳng song song thành hai đường thẳng song song hoặc trùng nhau.
  \item Bảo toàn tỉ số độ dài của hai đoạn thẳng nằm trên hai đường thẳng song song hoặc cùng nằm trên một đường thẳng: $\boxed{\dfrac{A'B'}{C'D'} = \dfrac{AB}{CD}}$.
  \end{itemize}
\item \textbf{Hình biểu diễn của một hình không gian:} Là hình chiếu song song của hình đó trên một mặt phẳng hoặc hình đồng dạng với hình chiếu đó. Quy tắc biểu diễn một số hình phẳng:
  \begin{itemize}
  \item Tam giác bất kì (đều, cân, vuông, thường) $\to$ tam giác bất kì.
  \item Hình vuông, hình chữ nhật, hình thoi, hình bình hành $\to$ hình bình hành.
  \item Hình thang $\to$ hình thang có tỉ số độ dài hai đáy không đổi.
  \item Đường tròn $\to$ đường elip hoặc đường tròn.
  \end{itemize}
\end{itemize}

\textbf{2. Phương pháp giải các dạng}
%% DANG 1H4?6-1
\textbf{Dạng 1. Câu hỏi lý thuyết}
\begin{enumerate}
\item Đọc kĩ giả thiết để xác định đối tượng được hỏi (định nghĩa, tính chất bảo toàn, vị trí tương đối giữa các hình chiếu).
\item Đối chiếu với 4 tính chất của phép chiếu song song (thẳng hàng, song song hoặc trùng nhau, tỉ số đoạn thẳng cùng phương).
\item Kiểm tra trường hợp đặc biệt: đường thẳng song song với phương chiếu có hình chiếu là một điểm; hai đường thẳng chéo nhau có hình chiếu có thể song song, cắt nhau hoặc trùng nhau.
\end{enumerate}
Lưu ý: Phép chiếu song song không bảo toàn độ dài đoạn thẳng, số đo góc và không bảo toàn tỉ số của hai đoạn thẳng không cùng phương.

%% DANG 1H4?6-2
\textbf{Dạng 2. Hình biểu diễn của một hình không gian}
\begin{enumerate}
\item Xác định các yếu tố bất biến của hình gốc: quan hệ song song, trung điểm, trọng tâm, tỉ số độ dài các đoạn thẳng cùng nằm trên một đường thẳng hoặc trên hai đường thẳng song song.
\item Chọn hình biểu diễn phù hợp theo quy tắc (hình vuông/chữ nhật/thoi $\to$ hình bình hành; tam giác $\to$ tam giác; hình thang $\to$ hình thang giữ nguyên tỉ số hai đáy).
\item Vẽ hình đảm bảo các yếu tố bất biến; dùng nét liền cho đường nhìn thấy và nét đứt cho đường bị che khuất.
\end{enumerate}
Lưu ý: Hình bình hành có thể biểu diễn hình vuông, hình chữ nhật, hình thoi nhưng không thể biểu diễn hình thang có hai đáy không bằng nhau.

%% DANG 1H4?6-3
\textbf{Dạng 3. Xác định yếu tố song song}
\begin{enumerate}
\item Xác định rõ mặt phẳng chiếu $(\alpha)$, phương chiếu $\Delta$ và điểm (hoặc đoạn thẳng, hình phẳng) cần tìm hình chiếu.
\item Qua mỗi điểm $M$ cần chiếu, dựng đường thẳng $d$ song song (hoặc trùng) với phương chiếu $\Delta$.
\item Tìm giao điểm $M' = d \cap (\alpha)$ bằng cách chọn mặt phẳng phụ chứa $d$ để tìm giao tuyến với $(\alpha)$, hoặc dùng tính chất bảo toàn trung điểm, song song, tỉ số để suy ra trực tiếp ảnh.
\end{enumerate}
Lưu ý: Nếu $M \in (\alpha)$ thì hình chiếu của $M$ là chính $M$. Nếu đoạn thẳng song song với phương chiếu $\Delta$ thì hình chiếu của nó chỉ là một điểm.

%% DANG 1H4?6-4
\textbf{Dạng 4. Xác định phương chiếu}
\begin{enumerate}
\item Tìm một điểm $M \notin (\alpha)$ và hình chiếu tương ứng $M' \in (\alpha)$ của nó qua phép chiếu song song đã cho.
\item Đường thẳng $MM'$ (hoặc mọi đường thẳng song song với $MM'$) chính là phương chiếu cần tìm.
\item Với bài toán có điều kiện đặc biệt (hai điểm $A, B$ có hình chiếu trùng nhau; hình chiếu của đa giác phẳng thuộc $(\beta)$ là một đoạn thẳng), suy ra phương chiếu $\Delta \parallel AB$ (hoặc $\Delta \parallel (\beta)$ hay $\Delta \subset (\beta)$).
\end{enumerate}
Lưu ý: Phương chiếu $\Delta$ bắt buộc phải cắt mặt phẳng chiếu $(\alpha)$, không được song song hoặc nằm trong $(\alpha)$.

%% DANG 1H4?6-5
\textbf{Dạng 5. Tính tỉ số đoạn thẳng, diện tích qua phép chiếu}
\begin{enumerate}
\item Xác định hình chiếu song song của các điểm, đoạn thẳng hoặc đa giác cần tính tỉ số lên mặt phẳng chiếu $(\alpha)$.
\item Đưa về các cặp đoạn thẳng cùng nằm trên một đường thẳng hoặc nằm trên hai đường thẳng song song, sau đó áp dụng tính chất bảo toàn tỉ số $\dfrac{A'B'}{C'D'} = \dfrac{AB}{CD}$ kết hợp định lí Ta-lét, tính chất trung điểm, trọng tâm.
\item Với tỉ số diện tích của hai hình phẳng cùng nằm trong một mặt phẳng (hoặc hai mặt phẳng song song), dùng tính chất tỉ số diện tích được bảo toàn qua phép chiếu song song: $\boxed{\dfrac{S_{\mathcal{H}_1'}}{S_{\mathcal{H}_2'}} = \dfrac{S_{\mathcal{H}_1}}{S_{\mathcal{H}_2}}}$.
\end{enumerate}
Lưu ý: Chỉ được áp dụng trực tiếp tỉ số đoạn thẳng khi hai đoạn thẳng gốc cùng phương (song song hoặc thẳng hàng); với hai đoạn thẳng cắt nhau hoặc chéo nhau, phải chiếu trung gian hoặc tính trực tiếp trên hình chiếu.
